\chapter{William Timothy Gowers}\label{ch1998williamtimothygowers:william-timothy-gowers-fields-medal-1998-a-comprehensive-report-on-his-work-and-impact}

\begin{center}\rule{0.5\linewidth}{0.5pt}\end{center}

\textbf{Prepared:} July 2026
\textbf{Subject:} William Timothy Gowers (born November 20, 1963, Marlborough, England)
\textbf{Position:} Professor of Mathematics, University of Cambridge (since 2003)
\textbf{Doctoral Advisor:} Bela Bollobas (University of Cambridge)
\textbf{Fields Medal Citation:} \emph{``For his contributions to functional
analysis and combinatorics.''}

\begin{center}\rule{0.5\linewidth}{0.5pt}\end{center}

\begin{prereqs}
Functional analysis (Banach spaces), combinatorics, additive
combinatorics, Szemeredi's theorem.
\end{prereqs}

\begin{objectives}
State the Gowers norms \cite{gowers2001} and their role in additive combinatorics;
explain Szemeredi's theorem and the density Hales--Jewett theorem;
describe the proof approach.
\end{objectives}

% ============================================================
\section{Problem History}
% ============================================================

\textbf{Szemeredi's theorem} (1975) states that every subset of
$\mathbb{Z}$ of positive upper density contains arbitrarily long
arithmetic progressions.  This landmark result was proved after
earlier work by Erdos and Turan and van der Waerden.

The \textbf{density Hales--Jewett theorem} (DHJ) is a
multidimensional generalization: for every $\delta>0$ and every
$k$, there exists $N$ such that any subset of $[k]^N$ of density
at least $\delta$ contains a \textbf{combinatorial line}---a set
$\{(x_1,\dots,x_N): x_i\in\{a_i,*\}\}$ where at least one
coordinate is a wildcard.

The original DHJ was proved by density increment arguments, but the
bounds were extremely weak:
$N\approx k^{2^{2^{2^{k}}}}$, a tower of exponentials.  This is
far too large for applications including quantitative versions of
Szemeredi's theorem.

The \textbf{Gowers norms} were introduced by Gowers in his 2001
proof of Szemeredi's theorem with quantitative bounds.  For
$f\colon\mathbb{Z}_N\to\mathbb{R}$:
\[
  \|f\|_{U^k}
  = \Biggl(\frac{1}{N^{k+1}}
    \sum_{x,h_1,\dots,h_k}
    \prod_{\omega\in\{0,1\}^k}
      \mathcal{C}^{|\omega|}
        f(x+\omega\cdot h)\Biggr)^{2^{-k}}
\]

Gowers introduced these norms while working on the
\textbf{Banach--Saks problem}, which asks whether every bounded
sequence in a Banach space has a subsequence with convergent Cesaro
mean.  He disproved this for certain Banach spaces using
Szemeredi-type arguments.

The connection between Szemeredi's theorem and functional analysis
was a surprise.  Gowers's innovation was to use Gowers norms as a
measure of uniformity, detecting whether a set behaves randomly
(small norm) or is structured (large norm).

The \textbf{Erdos conjecture} on arithmetic progressions states that
if $\sum_{a\in A}1/a=\infty$ for a set $A\subset\mathbb{N}$, then
$A$ contains arbitrarily long arithmetic progressions.  This was
proved by Green and Tao (2008) using Gowers norms and the inverse
conjecture.

The \textbf{Szemeredi regularity lemma} states that every large
graph can be approximated by a union of random-like bipartite
graphs.  Gowers showed that this lemma is essentially best possible
by constructing graphs where the number of parts must be a tower of
exponentials.  This led to the theory of \textbf{hypergraph
regularity}, which Gowers developed.

% ============================================================
\section{How Was It Solved}
% ============================================================

The proof of DHJ proceeds by \textbf{density increment} and
\textbf{energy increment}.  Given
$A\subseteq[k]^N$ of density $\alpha$, the goal is to find a
combinatorial line or a large subspace with significantly higher
density.

The key tool is the Gowers norm $U^2$, related to the $L^4$ norm of
the Fourier transform.  If $A$ has small $U^2$ norm, it behaves
like a random set and combinatorial lines can be found by counting.
If $A$ has large $U^2$ norm, it has strong Fourier bias and density
can be increased on a subspace.

The \textbf{density increment} step: if $A$ has large $U^2$ norm,
there exists a large subspace $W\subseteq[k]^N$ on which density
increases from $\alpha$ to $\alpha+\epsilon$.

The \textbf{energy increment} step: if density cannot be increased,
the $U^2$ norm is small and a combinatorial line is found directly.

The proof iterates these steps.  After $O(1/\epsilon)$ iterations,
either density reaches 1 or the $U^2$ norm becomes small.  The bound
obtained is $N\approx k^{Ck^2}$, much smaller than the original
tower-of-exponentials.

The \textbf{Gowers norm} $U^k$ can be defined inductively.  For
$k=1$, $\|f\|_{U^1}=|\mathbb{E}_x f(x)|$.  For $k\geq 2$:
\[
  \|f\|_{U^k}^{2^k}
  = \mathbb{E}_x\,\mathbb{E}_{h_1}\cdots\mathbb{E}_{h_k}
    \prod_{\omega\in\{0,1\}^k}
      \mathcal{C}^{|\omega|}f(x+\omega\cdot h)
\]
where $\mathcal{C}$ denotes complex conjugation.  The $U^2$ norm
satisfies
$\|f\|_{U^2}^4
 = \mathbb{E}_{x,h}|f(x)\overline{f(x+h)}|^2$, which is closely
related to the $L^4$ norm of the Fourier transform.

The \textbf{inverse theorem for Gowers norms} characterizes
functions with large Gowers norm.  For $U^2$, a large norm implies
a large Fourier coefficient.  For $U^3$, large norm implies
correlation with a \textbf{quadratic phase}
$e(ax^2+bx+c)$.  In general, functions with large $U^k$ norm
correlate with \textbf{nilsequences}---functions defined on $k-1$
step nilpotent groups.

% ============================================================
\section{Impact of This Discovery}
% ============================================================

The density Hales--Jewett theorem with quantitative bounds has
become a cornerstone of additive combinatorics.  The bounds were
improved to $N\approx k^{Ck}$ by the polymath project, one of the
first major applications of massively collaborative mathematics.

Gowers norms have become fundamental tools in additive
combinatorics, connecting Fourier analysis, ergodic theory, and
combinatorics.  They detect structure in sets and functions, used in
sum-product phenomena and arithmetic regularity.

Gowers's work initiated modern additive combinatorics.  The
Green--Tao theorem on primes in arithmetic progressions (2004) uses
Gowers norms as a central tool, and the theory of quadratic
uniformity is based on Gowers norm analysis.

The connection between Banach spaces and combinatorics has led to new
results in both areas.  Gowers spaces $U^k$ provide a bridge between
Banach space geometry and arithmetic set structure.

Beyond specific results, Gowers demonstrated the power of
quantitative methods in combinatorics, replacing ad hoc arguments
with a unified theory.

The \textbf{Gowers--Szemeredi theorem} (1990), proved before the
norm paper, showed that every set with no three-term arithmetic
progression has sumset growth at least
$|A+A|\geq |A|^{1+\epsilon}$.  This was a precursor to the full
Gowers norm framework.

Gowers's \textbf{hypergraph regularity lemma} provides a
generalization of the Szemeredi regularity lemma to $k$-uniform
hypergraphs, with applications to the generalized Szemeredi theorem.

The \textbf{Balog--Szemeredi--Gowers lemma} states that if
$|A\cdot A|$ and $|A+A|$ are both close to $|A|$, then $A$ has a
large subset with small sumset.  This lemma is fundamental in
additive combinatorics.

% ============================================================
\section{Modern Views and Improvements}
% ============================================================

The \textbf{polymath project}, initiated by Gowers in 2009, used
massively collaborative research to improve DHJ bounds to
$N\approx k^{Ck}$ through a combination of Gowers's arguments, the
triangle removal lemma, and new combinatorial techniques.

The \textbf{inverse theorem} for Gowers norms, characterizing
functions with large Gowers norm, has been proved in several cases
by Green, Tao, Ziegler, and others.  The full inverse theorem for
$U^k$ remains open.

The \textbf{Green--Tao theorem} uses the Gowers norm framework as a
central tool.  The proof decomposes primes into structured,
quadratic, and nilsequence parts.

The \textbf{Polymath8} project improved bounds on the smallest
prime in arithmetic progressions of length $k$, demonstrating the
power of collaborative approaches.

Recent work by Sanders, Tao, and others on the \textbf{inverse
theorem for $U^3$} has shown that functions with large $U^3$ norm
correlate with cubic phases, completing the inverse theorem for
$k=3$.  The general case for $k\geq 4$ remains one of the central
open problems in additive combinatorics.

Gowers continues contributing to additive combinatorics, including
inverse theorems, approximate group structure, and prime
distribution.  His 2012 paper with Tao on approximate groups uses
Gowers norms and Bourgain--Chang results.

The Erdos--Szemeredi conjecture remains open, but Gowers norm
methods have made significant progress on related problems.

% ============================================================
\begin{center}\rule{0.5\linewidth}{0.5pt}\end{center}

\section*{Summary}
\addcontentsline{toc}{section}{Summary}

Gowers introduced Gowers norms and proved quantitative density
Hales--Jewett and Szemeredi theorems, connecting Banach space
theory to combinatorics.  His work initiated modern additive
combinatorics and provided tools used in the Green--Tao theorem and
the polymath project.

% ============================================================
